\renewcommand{\Anno}{2010}

\renewcommand{\NumeroProblema}{3}

\renewcommand{\Testo}{ 
Paolo e Francesca stanno giocando a testa o croce con una moneta.  Se
esce testa, Paolo guadagna un punto, se esce croce, Francesca guadagna
un punto.  La partita prosegue fino a quando uno dei due giocatori
arriva a 6 punti, a meno che l'altro giocatore non abbia gi\`{a} 5
punti, nel qual caso si prosegue fino a quando il punteggio per la
prima volta diventa 7--5 o 7--6 (come accade nei set del tennis, per
intendersi).
\begin{enumerate}
	\item  Determinare la probabilit\`{a} che la partita termini con 
	il punteggio di 6--1 (indipendentemente da chi vince).
	
	\item Determinare la probabilit\`{a} che durante la partita si
	realizzi il punteggio 3--3.

	\item Sempre indipendentemente dal vincitore, indichiamo con $P$ la
	probabilit\`{a} che la partita si concluda sul 6--4, con $Q$ la
	probabilit\`{a} che si concluda sul 7--5, e con $R$ la
	probabilit\`{a} che si concluda sul 7--6.
	
	Determinare l'ordinamento tra i numeri $P$, $Q$, $R$.
\end{enumerate}
}

\renewcommand{\Risposta}{
\begin{enumerate}
	\renewcommand{\labelenumi}{(\alph{enumi})} 
	\item  3/32

	\item  5/16
	
	\item  $P>Q=R$, e anzi $\displaystyle{Q=R=\frac{P}{2}}$
\end{enumerate}
}

\renewcommand{\Aiutino}{
\begin{enumerate}
	\renewcommand{\labelenumi}{(\alph{enumi})} 
	\item Quanti sono i casi possibili? Quanti sono i casi 
	favorevoli? Attenzione: l'ultimo punto \`{e} stato sicuramente 
	realizzato dal vincitore.

	\item Quanti sono i casi possibili? Quanti sono i casi 
	favorevoli? Quelli che contano sono solo i primi 6 lanci \ldots 

	\item Segnaliamo due approcci.
	\begin{itemize}
		\item Calcolare esplicitamente $P$, $Q$, $R$ mediante casi
		possibili e casi favorevoli.  Bisogna per\`{o} stare attenti a
		quali posizioni possono occupare, nella sequenza dei lanci,
		quelli favorevoli allo sconfitto.  Ad esempio, in caso di
		finale 7--5, in quali posizioni possono trovarsi i lanci
		favorevoli allo sconfitto?
	
		\item Esprimere $P$, $Q$, $R$ in funzione di $p_{5-4}$,
		cio\`{e} in funzione della probabilit\`{a} che la partita
		passi sul punteggio di 5--4 (ci deve passare per forza?).
		Esaminando i possibili andamenti dei lanci successivi, si
		riesce facilmente a dimostrare (come?)  che 
		$$P=\frac{1}{2}\,p_{5-4}
		\quad\mbox{e}\quad
		Q=R=\frac{1}{4}\,p_{5-4}.$$
	\end{itemize}

\end{enumerate}
}

\renewcommand{\SoluzioneUfficiale}{\TitoloSoluzione\DomandaA 
La probabilit\`{a} che la partita termini con il punteggio di 6--1 
\`{e} 3/32.

Calcoliamo intanto la probabilit\`{a} che la partita finisca 6--1 a
favore di Paolo.  Questo punteggio finale vuol dire che nei primi 7
lanci ci sono state 6 testa ed una croce.  Abbiamo quindi $2^{7}=128$
sequenze possibili, e 6 sequenze favorevoli, a seconda del lancio in
cui \`{e} uscita la croce (che pu\`{o} essere uno qualunque dei primi
6 lanci, ma non l'ultimo, altrimenti Paolo avrebbe gi\`{a} vinto
6--0).

La probabilit\`{a} che Paolo vinca 6--1 \`{e} quindi $6/128=3/64$. 
Poich\'{e} lo stesso discorso vale per Francesca, non resta che 
moltiplicare per 2.

\DomandaB 

Chiaramente la partita pu\`{o} trovarsi sul 3--3 solo dopo 6 lanci.  I
casi possibili sono dunque $2^{6}=64$.  I casi favorevoli sono tanti
quanti i modi di scegliere i 3 lanci su 6 in cui esce, per
esempio, testa, e sono quindi $\binom{6}{3}=20$.

Di conseguenza il punteggio passa per il 3--3 con probabilit\`{a} 
$20/64=5/16$.

\DomandaC 
Mostriamo che
\begin{equation}
	Q=R=\frac{1}{2}P,
	\label{PQR}
\end{equation}
dunque in particolare $P>Q=R$.

Per prima cosa osserviamo che le partite che terminano sul 6--4, 7--5 o
7--6 passano tutte per il punteggio di 5--4 (a favore di uno o
dell'altro).  A tal fine basta considerare il primo momento in cui
qualcuno arriva a 5: se costui passa a 6 prima che l'altro sia
arrivato a 4, evidentemente la partita finisce su un altro punteggio.

Indichiamo ora con $p_{5-4}$ la probabilit\`{a} che la partita passi
sul 5--4, ed esaminiamo il lancio successivo (il decimo).  Se \`{e}
favorevole al giocatore che sta a 5, la partita finisce sul 6--4 per
lui, quindi
\begin{equation}
	P=\frac{1}{2}p_{5-4}.
	\label{prob-P}
\end{equation}

Se il decimo lancio \`{e} favorevole al giocatore che sta a 4, allora 
la partita va sul 5--5, e poi sul 6--5 al lancio successivo. Con 
ovvio significato dei simboli, avremo pertanto che
$$p_{6-5}=p_{5-5}=\frac{1}{2}p_{5-4}.$$

A questo punto il dodicesimo lancio \`{e} decisivo: se \`{e} a favore 
di chi sta in vantaggio, la partita termina sul 7--5, altrimenti si 
va sul 6--6 ed al passaggio successivo finir\`{a} comunque 7--6 per 
qualcuno. Ne segue che
\begin{equation}
	Q=R=\frac{1}{2}p_{6-5}=\frac{1}{2}p_{5-5}=\frac{1}{4}p_{5-4}.
	\label{prob-QR}
\end{equation}

Dalla (\ref{prob-P}) e dalla (\ref{prob-QR}) segue immediatamente la 
(\ref{PQR}).
}

\renewcommand{\SoluzioniAlternative}{\TitoloSoluzioneAlternativaA

\`{E} ovviamente possibile rispondere alla domanda~(c) calcolando 
esplicitamente le probabilit\`{a} $P$, $Q$, $R$. Un modo di farlo 
\`{e} di determinare quella che prima abbiamo chiamato $p_{5-4}$, 
e poi procedere come nella soluzione precedente. In alternativa, si 
pu\`{o} fare il calcolo direttamente, prestando particolare attenzione 
a come si contano i casi favorevoli.

Partiamo da $P$.  I casi possibili sono $2^{10}=1024$.  I casi
favorevoli sono tanti quanti i modi di scegliere i 4 lanci che hanno
portato i punti allo sconfitto tra le nove posizioni possibili
(l'ultimo lancio \`{e} infatti sicuramente favorevole al vincitore),
moltiplicati per 2 perch\'{e} lo sconfitto pu\`{o} essere Paolo o
Francesca.  Otteniamo cos\`{i} che
$$P=\frac{1}{1024}\cdot\binom{9}{4}\cdot 2=\frac{63}{256}.$$

Passiamo ora a $Q$.  I casi possibili sono $2^{12}=4096$.  Il computo
dei casi favorevoli richiede particolare attenzione.  Dobbiamo
scegliere i 5 lanci favorevoli allo sconfitto, ma nel fare questo
bisogna tenere conto che dopo 10 lanci la partita stava 5--5
(altrimenti sarebbe gi\`{a} finita), dunque i 5 lanci vanno scelti
nelle prime 10 posizioni (in qualunque modo, perch\'{e} non importa
come si \`{e} arrivati sul 5--5).  Moltiplicando come prima per 2,
otteniamo cos\`{i} che
$$Q=\frac{1}{4096}\cdot\binom{10}{5}\cdot 2=\frac{63}{512}.$$

Concludiamo infine con $R$. I casi possibili sono $2^{13}=8192$. Come 
nel caso precedente i casi possibili sono tutti quelli che portano 
sul 5--5, da moltiplicare per 2 a seconda del fatto che l'undicesimo 
lancio sia stato a favore del vincitore o dello sconfitto (l'esito 
dei lanci dal dodicesimo in poi \`{e} obbligato), e poi ancora per 2 
perch\'{e} i possibili vincitori sono 2. Otteniamo cos\`{i} che
$$R=\frac{1}{8192}\cdot\binom{10}{5}\cdot 2=\frac{63}{512}.$$

Ovviamente le 3 probabilit\`{a} trovate soddisfano le relazione 
(\ref{PQR}).
}

\renewcommand{\Osservazioni}{\TitoloOsservazione

\begin{itemize}
	\item La seguente tabella riassume le probabilit\`{a} di tutti i
	possibili punteggi finali, espresse in 512-esimi per poter essere
	confrontate pi\`{u} facilmente.  Si vede, tra l'altro, che il 6--4
	\`{e} di poco il punteggio pi\`{u} probabile, seguito a ruota dal 
	6--3.
	\begin{center}
		\renewcommand{\arraystretch}{1.5}
		\begin{tabular}{|c||c|c|c|c|c|c|c|}
			\hline
			Punteggio & 6--0 & 6--1 & 6--2 & 6--3 & 6--4 & 7--5 & 7--6  \\
			\hline
			Probabilit\`{a} & 16/512 & 48/512 & 84/512 & 112/512 &
			126/512 & 63/512 & 63/512 \\
			\hline
		\end{tabular}
	\end{center}

	\item La soluzione riportata al punto~(c) funziona allo stesso
	modo anche se si rimpiazza il numero 6 con un qualunque intero
	$n\geq 2$, cio\`{e} se la partita viene vinta dal primo giocatore
	che arriva ad $n$ con almeno due punti di distacco, oppure dal
	primo giocatore che arriva ad $(n+1)$.  Ovviamente in questo caso
	i punteggi che definiscono le probabilit\`{a} $P$, $Q$, $R$
	sarebbero $n$--$(n-2)$, $(n+1)$--$(n-1)$ e $(n+1)$--$n$.
	\end{itemize}
}

\renewcommand{\Commenti}{
}
